\section*{Overview}

These notes are the part of the competitive programming archive that is not tied to one contest year.
The goal is simple: keep the ideas that I actually want to reuse in a form that is still readable after
the adrenaline of the original contest is gone.

I care about three things here:

\begin{itemize}
  \item the structural reason an algorithm or data structure works,
  \item the implementation shape I would still trust in contest code,
  \item the failure modes that usually cause wrong answers or wasted time.
\end{itemize}

\section*{How the Notes Are Written}

Each topic keeps \texttt{note.tex} as the written source of truth and \texttt{code.cpp} as the implementation
source of truth. The site renders the TeX into HTML and shows the C++ file separately, so the prose can stay
focused while the code remains directly copyable.

The finished notes are meant to be public-facing study material, not private scraps. That means the emphasis is
on explanations, invariants, use cases, and pitfalls instead of dumping templates with no context.

\section*{Recommended Learning Path}

If I were revising this section from scratch, I would split it into a few routes instead of pretending there is one
universal order:

\begin{itemize}
  \item \textbf{core first}: binary search, prefix sums, sliding window, coordinate compression, Fenwick tree, BFS / DFS,
  \item \textbf{graph path}: topological sort, Dijkstra, MST, SCC, Euler tour, LCA, tree decompositions,
  \item \textbf{DP optimization path}: knapsack, tree DP, monotone queue, divide-and-conquer DP, Knuth, convex hull trick, rerooting,
  \item \textbf{advanced data structures}: segment tree, sparse table, ordered set, Li Chao tree, persistent segment tree, rollback DSU,
  \item \textbf{number theory path}: extended Euclid, modular arithmetic, CRT, sieve, Mobius, Miller-Rabin, discrete log.
\end{itemize}

That is still not the only valid order, but it is closer to how I actually revise: pick a branch, stay inside it long
enough for the invariants to become familiar, then switch.

\section*{What ``Finished'' Means Here}

A finished note should answer the questions I usually care about when revising:

\begin{itemize}
  \item What kind of contest situation makes this the right tool?
  \item What invariant is doing the real work?
  \item Which implementation shape is the least fragile under time pressure?
  \item Which edge cases usually break the first version?
  \item What does a worked contest-style example actually look like?
\end{itemize}

Some topics are still only planned. They are listed openly as planned instead of pretending that a shallow page
is enough.

\section*{References Behind the Section}

The explanations are rewritten in my own words, but the topic map and many implementation choices were informed by:

\begin{itemize}
  \item \href{https://codeforces.com/catalog}{Codeforces Catalog},
  \item \href{https://cp-algorithms.com/}{cp-algorithms},
  \item \href{https://usaco.guide/}{USACO Guide},
  \item \href{https://usaco.guide/CPH.pdf}{Competitive Programmer's Handbook},
  \item \href{https://github.com/kth-competitive-programming/kactl}{KACTL},
  \item \href{https://oi-wiki.org/}{OI Wiki}.
\end{itemize}
